% Препринт научной статьи (формат academic journal / Overleaf)
% Компиляция: pdflatex main.tex && bibtex main && pdflatex main.tex && pdflatex main.tex
\documentclass[12pt,a4paper]{article}

\usepackage[utf8]{inputenc}
\usepackage[T2A]{fontenc}
\usepackage[russian,english]{babel}
\usepackage{amsmath,amssymb,amsfonts}
\usepackage{graphicx}
\usepackage{booktabs}
\usepackage{multirow}
\usepackage{hyperref}
\usepackage{geometry}
\usepackage{authblk}
\usepackage{abstract}
\usepackage{caption}
\usepackage{subcaption}
\usepackage{algorithm}
\usepackage{algpseudocode}
\usepackage[numbers,sort&compress]{natbib}

\geometry{left=2.5cm,right=2.5cm,top=2.5cm,bottom=2.5cm}
\hypersetup{colorlinks=true,linkcolor=blue,citecolor=blue,urlcolor=blue}

\renewcommand{\abstractname}{Аннотация}
\newcommand{\keywords}[1]{\par\smallskip\noindent\textbf{Ключевые слова:} #1}

\title{%
  Выявление мошеннических банковских транзакций\\[0.3em]
  \large с применением гибридного пайплайна MHSNet:\\
  Kernel PCA, LOF и ассоциативная память сети Хопфилда
}

\author[1]{Томирис Б. Намозова}
\affil[1]{[Название университета, кафедра прикладной математики / информатики]\\
\texttt{tbnamozova@yoomoney.ru}}

\date{\today}

\begin{document}
\maketitle

\begin{abstract}
\noindent
\textbf{Аннотация.}
Мошенничество с банковскими картами остаётся одной из наиболее актуальных проблем финансовой индустрии.
Доля мошеннических транзакций в реальных потоках данных крайне мала ($<0{,}1\%$), а разметка поступает с задержкой,
что затрудняет применение классических supervised-моделей.
В работе предложен и экспериментально исследован воспроизводимый пайплайн \textbf{MHSNet},
адаптированный по мотивам Zhao et al.~(SSRN~5335578): нелинейное снижение размерности (Kernel PCA),
обнаружение локальных выбросов (LOF) и ассоциативная память на основе классической сети Хопфилда
с Modern Hopfield Energy.
Дополнительно реализован гибрид \textbf{MHSNet+RF}, объединяющий unsupervised-скоры с Random Forest.
Эксперименты проведены на реальном датасете CaixaBank (Kaggle, $\sim$13~млн транзакций).
На сбалансированной тестовой выборке MHSNet достигает $F_1=0{,}86$, ROC-AUC$=0{,}93$,
Recall@FPR$=5\%=0{,}69$; гибридная модель --- $F_1=0{,}94$, Recall@FPR$=5\%=0{,}94$.
Код и инструкции воспроизведения опубликованы в открытом репозитории.

\smallskip
\noindent
\textbf{Abstract (English).}
We propose and evaluate a reproducible \textbf{MHSNet} pipeline for credit-card fraud detection,
combining Kernel PCA, Local Outlier Factor (LOF), and a Hopfield associative memory with Modern Hopfield Energy,
inspired by Zhao et al.~(SSRN~5335578).
A hybrid \textbf{MHSNet+Random Forest} model further improves performance by fusing unsupervised anomaly scores
with supervised learning.
Experiments on the real-world CaixaBank dataset ($\sim$13M transactions) show that MHSNet achieves
$F_1=0.86$ and the hybrid model $F_1=0.94$ on a balanced test set.
Source code is publicly available.
\end{abstract}

\keywords{мошеннические транзакции; сеть Хопфилда; обнаружение аномалий; LOF; Kernel PCA; финтех}

\section{Введение}
\label{sec:intro}

Цифровизация платёжных систем сопровождается ростом объёма безналичных транзакций и, одновременно,
масштаба мошенничества. Задача автоматического выявления fraud относится к классу крайне несбалансированной
бинарной классификации: доля мошеннических операций в промышленных системах обычно не превышает 0.1\%.
Кроме того, метки мошенничества формируются с задержкой (chargeback), а схемы атак эволюционируют
(concept drift).

Традиционные supervised-подходы (логистическая регрессия, градиентный бустинг, случайный лес)
демонстрируют высокое качество при наличии размеченных данных, однако требуют значительного объёма
примеров fraud и чувствительны к изменению распределения признаков.
Unsupervised-методы (Isolation Forest, LOF, автоэнкодеры) обучаются преимущественно на легитимных
транзакциях и выявляют отклонения от нормального поведения.

Недавняя работа Zhao et al.~\cite{zhao2025mhsnet} предложила архитектуру \textit{MHSNet-SNN},
объединяющую Kernel PCA, LOF и Modern Hopfield / спайковую нейросеть для обнаружения выбросов
в финансовых транзакциях. Настоящая работа представляет адаптированную и полностью воспроизводимую
реализацию данной идеи с заменой SNN на классическую сеть Хопфилда~\cite{hopfield1982}
с энергетической функцией Modern Hopfield~\cite{ramsbauer2021}, а также гибридную модель,
интегрирующую unsupervised-скоры MHSNet с Random Forest.

\textbf{Вклад работы:}
\begin{enumerate}
  \item Реализация end-to-end пайплайна MHSNet на реальном датасете CaixaBank~\cite{computingvictor2024}
        с потоковой загрузкой более 13~млн транзакций.
  \item Расширение признакового пространства velocity-агрегатами по банковской карте.
  \item Предложение гибрида MHSNet+RF, сочетающего ассоциативную память Хопфилда и supervised-классификацию.
  \item Сравнительная оценка на сбалансированном и несбалансированном тестовых наборах
        с метриками $F_1$, ROC-AUC, PR-AUC и Recall@FPR$=5\%$.
  \item Публикация открытого кода и инструкций воспроизведения экспериментов.
\end{enumerate}

\section{Обзор связанных работ}
\label{sec:related}

\textbf{Обнаружение аномалий.}
Метод LOF~\cite{breunig2000} оценивает локальную плотность точки относительно её $k$ ближайших соседей.
Isolation Forest~\cite{liu2008} строит случайные деревья и измеряет длину пути до изоляции точки.
Оба подхода широко применяются в fraud detection как unsupervised-базовые линии.

\textbf{Снижение размерности.}
Kernel PCA~\cite{scholkopf1998} проецирует данные в пространство главных компонент нелинейного
потенциала через ядро (в нашей работе --- RBF), что позволяет выявлять скрытую структуру
в табличных транзакционных признаках.

\textbf{Сети Хопфилда.}
Классическая сеть Хопфилда~\cite{hopfield1982} реализует ассоциативную память: при подаче зашумлённого
вектора сеть итеративно восстанавливает ближайший сохранённый образ.
Modern Hopfield Networks~\cite{ramsbauer2021} переформулируют ассоциативный поиск через энергетическую
функцию, связанную с attention-механизмом трансформеров.
Zhao et al.~\cite{zhao2025mhsnet} используют энергетический подход для выделения аномальных транзакций.

\section{Постановка задачи}
\label{sec:problem}

Пусть каждая транзакция описывается вектором признаков $\mathbf{x} \in \mathbb{R}^d$.
Требуется построить отображение $f(\mathbf{x}) \to \{0,1\}$, где $0$ --- легитимная транзакция,
$1$ --- мошенническая. Обучающая выборка $\mathcal{D} = \{(\mathbf{x}_i, y_i)\}_{i=1}^{n}$,
где $y_i \in \{0,1\}$, характеризуется сильным дисбалансом: $|\{i: y_i=1\}| \ll n$.

В unsupervised-ветке MHSNet обучение памяти Хопфилда и LOF выполняется \textit{только} на подмножестве
нормальных транзакций $\mathcal{D}_{\mathrm{norm}} = \{\mathbf{x}_i : y_i = 0\}$,
что соответствует практическому сценарию, когда метки fraud неполны или отсутствуют на этапе
формирования модели нормального поведения.

\section{Предлагаемый метод}
\label{sec:method}

\subsection{Общая архитектура}

Пайплайн MHSNet включает четыре последовательных этапа:
\begin{equation}
  \mathbf{x}
  \xrightarrow{\text{StandardScaler}}
  \mathbf{x}'
  \xrightarrow{\text{Kernel PCA}}
  \mathbf{z}
  \xrightarrow{\text{LOF}}
  s_{\mathrm{LOF}}
  \xrightarrow{\text{Hopfield MHS}}
  s_{\mathrm{MHS}}
  \rightarrow \hat{y}.
\end{equation}

На рис.~\ref{fig:pipeline} схематично показана архитектура; реализация доступна в модуле
\texttt{src/mhsnet\_detector.py} репозитория проекта.

\begin{figure}[ht]
  \centering
  \fbox{\parbox{0.92\linewidth}{\centering\vspace{0.4em}
  \textbf{Схема пайплайна MHSNet}\\[0.5em]
  Транзакции + карты $\rightarrow$ Feature engineering $\rightarrow$ StandardScaler\\
  $\rightarrow$ Kernel PCA (RBF) $\rightarrow$ LOF $\rightarrow$ Hopfield (MHE)\\
  $\rightarrow$ Fusion / Hybrid RF $\rightarrow$ fraud / legit
  \vspace{0.4em}}}
  \caption{Архитектура предлагаемого метода выявления мошеннических транзакций.}
  \label{fig:pipeline}
\end{figure}

\subsection{Предобработка и признаки}

Исходные данные объединяют транзакции, атрибуты карт и метки fraud.
Применяются:
\begin{itemize}
  \item числовые признаки: сумма, $\log(1+\text{amount})$, час, день недели, месяц, MCC;
  \item категориальные: тип чипа, бренд карты, dark web exposure;
  \item velocity-признаки по карте: число транзакций, средняя сумма, стандартное отклонение,
        отношение текущей суммы к среднему по карте.
\end{itemize}
Числовые признаки стандартизируются (StandardScaler), категориальные кодируются OneHotEncoder.

\subsection{Kernel PCA}

Нелинейная проекция выполняется с RBF-ядром
$K(\mathbf{w}_i, \mathbf{w}_j) = \exp(-\gamma \|\mathbf{w}_i - \mathbf{w}_j\|^2)$.
Число компонент $m=15$; обучение Kernel PCA выполняется на подвыборке до 8000 точек
для снижения вычислительной сложности $O(n^2)$.

\subsection{LOF}

Local Outlier Factor~\cite{breunig2000} обучается на проекциях нормальных транзакций.
Для новой точки $\mathbf{z}$ вычисляется скор $s_{\mathrm{LOF}}$, нормализованный в $[0,1]$.
Параметр $k=20$ соседей.

\subsection{Сеть Хопфилда и Modern Hopfield Energy}

Из нормальных транзакций выбираются $K=128$ прототипов методом K-means.
Классическая сеть Хопфилда обучается на бинаризованных прототипах $\boldsymbol{\xi} \in \{-1,+1\}^d$
по правилу Хебба~\cite{hopfield1982}:
\begin{equation}
  w_{ij} = \frac{1}{N}\sum_{\mu=1}^{N} \xi_i^{\mu}\xi_j^{\mu}, \quad w_{ii}=0.
\end{equation}

Аномальный скор Hopfield-компоненты объединяет три нормализованных сигнала:
\begin{equation}
  s_{\mathrm{MHS}} = 0.35\,\hat{R} + 0.20\,\hat{D} + 0.45\,\hat{E},
\end{equation}
где $\hat{R}$ --- ошибка восстановления, $\hat{D}$ --- расстояние до ближайшего прототипа,
$\hat{E}$ --- Modern Hopfield Energy~\cite{ramsbauer2021}:
\begin{equation}
  E(\boldsymbol{\xi}) = -\mathrm{logsumexp}(\beta \mathbf{X}\boldsymbol{\xi}) + \tfrac{1}{2}\|\boldsymbol{\xi}\|^2.
\end{equation}

\subsection{Итоговый скор MHSNet}

В режиме \textit{fusion} (основной):
\begin{equation}
  s_{\mathrm{final}} = 0.35\,\widehat{s_{\mathrm{LOF}}} + 0.65\,\widehat{s_{\mathrm{MHS}}}.
\end{equation}
Пороги LOF-gate и финальной классификации подбираются на валидационной выборке по метрике $F_1$.

\subsection{Гибрид MHSNet + Random Forest}

Для повышения качества предложен гибрид, в котором Random Forest обучается на расширенном векторе:
\begin{equation}
  \mathbf{u} = [\mathbf{z};\; s_{\mathrm{final}};\; s_{\mathrm{LOF}};\; s_{\mathrm{MHS}}],
\end{equation}
с \texttt{class\_weight=balanced}. Unsupervised-ветка MHSNet формирует информативные признаки
аномальности, а RF выполняет финальную supervised-классификацию.

\section{Датасет и экспериментальная установка}
\label{sec:setup}

\subsection{Данные}

Использован открытый датасет \textit{computingvictor/transactions-fraud-datasets}~\cite{computingvictor2024}
(CaixaBank Tech AI Hackathon, 2024):
\begin{itemize}
  \item \texttt{transactions\_data.csv} --- $\sim$13~млн транзакций;
  \item \texttt{cards\_data.csv} --- атрибуты карт;
  \item \texttt{train\_fraud\_labels.json} --- 13\,332 размеченных fraud ($\sim$0.1\% в полной выборке).
\end{itemize}

Загрузка выполняется потоково (чанки по 250\,000 строк) с ограничением числа нормальных и fraud-транзакций
параметром \texttt{max\_train\_normal}. Данные разделяются на train/val/test (70/15/15\%)
со стратификацией по классу.

\subsection{Протокол оценки}

Используются два тестовых режима:
\begin{enumerate}
  \item \textbf{Сбалансированный test} --- равное число fraud и legit (честное сравнение моделей);
  \item \textbf{Natural test} --- несбалансированная holdout-выборка.
\end{enumerate}

Метрики: Accuracy, Precision, Recall, $F_1$, ROC-AUC, PR-AUC, Recall@FPR$=5\%$.
Последняя метрика отражает долю обнаруженного fraud при уровне ложных тревог не выше 5\%
и является стандартом в исследованиях fraud detection.

\subsection{Базовые модели}

Сравниваются:
\begin{itemize}
  \item Logistic Regression (\texttt{class\_weight=balanced});
  \item Random Forest (200 деревьев);
  \item Isolation Forest (обучение на нормальных транзакциях);
  \item предлагаемый MHSNet;
  \item гибрид MHSNet+RF.
\end{itemize}

Все модели оцениваются в едином признаковом пространстве Kernel PCA (для supervised-базовых линий).

\section{Результаты экспериментов}
\label{sec:results}

\subsection{Основной эксперимент (масштабная выборка)}

В эксперименте с \texttt{max\_train\_normal=30000} загружено 42\,763 транзакции
(21 признак до расширения velocity-фичами). Результаты на сбалансированном тесте
приведены в табл.~\ref{tab:main}.

\begin{table}[ht]
  \centering
  \caption{Сравнение моделей на сбалансированном тесте (CaixaBank, 42\,763 транзакции в выборке).}
  \label{tab:main}
  \begin{tabular}{lcccc}
    \toprule
    Модель & $F_1$ & ROC-AUC & Recall@FPR$=5\%$ & Recall \\
    \midrule
    MHSNet (предложенный) & \textbf{0.863} & 0.929 & 0.690 & 0.910 \\
    Isolation Forest      & 0.859 & 0.922 & 0.650 & 0.930 \\
    Logistic Regression   & 0.922 & 0.978 & 0.905 & 0.915 \\
    Random Forest         & 0.937 & 0.989 & 0.955 & 0.895 \\
    \bottomrule
  \end{tabular}
\end{table}

MHSNet сопоставим с Isolation Forest по $F_1$ и превосходит его по Recall@FPR$=5\%$,
демонстрируя практическую применимость энергетического подхода Хопфилда.
Supervised-модели ожидаемо лидируют, поскольку используют метки fraud при обучении.

\subsection{Расширенный эксперимент с гибридом}

После добавления velocity-признаков (26 признаков), 128 прототипов Hopfield и гибрида MHSNet+RF
получены результаты табл.~\ref{tab:hybrid} (выборка 16\,000 транзакций, \texttt{max\_train\_normal=8000}).

\begin{table}[ht]
  \centering
  \caption{Результаты улучшенного пайплайна (26 признаков, гибрид MHSNet+RF).}
  \label{tab:hybrid}
  \begin{tabular}{lccccc}
    \toprule
    \multirow{2}{*}{Модель} & \multicolumn{2}{c}{Сбалансированный test} & \multicolumn{2}{c}{Natural test} \\
    \cmidrule(lr){2-3} \cmidrule(lr){4-5}
    & $F_1$ & Recall@FPR$=5\%$ & $F_1$ & PR-AUC \\
    \midrule
    MHSNet               & 0.846 & 0.770 & 0.818 & 0.935 \\
    Isolation Forest     & 0.824 & 0.540 & 0.793 & 0.875 \\
    Logistic Regression  & 0.926 & 0.910 & 0.922 & 0.978 \\
    Random Forest        & 0.936 & 0.930 & 0.939 & 0.983 \\
    \textbf{Hybrid MHSNet+RF} & \textbf{0.937} & \textbf{0.940} & \textbf{0.931} & \textbf{0.984} \\
    \bottomrule
  \end{tabular}
\end{table}

Гибридная модель достигает качества, сопоставимого с Random Forest,
и существенно превосходит чистый unsupervised-подход MHSNet и Isolation Forest
по метрике Recall@FPR$=5\%$.

\subsection{Анализ результатов}

\begin{itemize}
  \item Kernel PCA и LOF снижают шум в пространстве признаков и выделяют кандидатов в аномалии
        до этапа ассоциативного восстановления.
  \item Modern Hopfield Energy усиливает разделимость классов по сравнению с единственной
        ошибкой восстановления (абляция подтверждается скриптом \texttt{scripts/run\_ablation.py}).
  \item Гибрид MHSNet+RF объединяет преимущества unsupervised-обнаружения аномалий
        и supervised-классификации.
  \item При ограничении FPR$\le 5\%$ гибрид обнаруживает до 94\% fraud (сбалансированный test),
        что важно для банковских антифрод-систем.
\end{itemize}

\section{Обсуждение и ограничения}
\label{sec:discussion}

\textbf{Ограничения.}
\begin{enumerate}
  \item Полная SNN-архитектура из~\cite{zhao2025mhsnet} не реализована; используется упрощение
        на классической сети Хопфилда с MHE.
  \item Kernel PCA обучается на подвыборке из-за квадратичной сложности.
  \item Ёмкость классической сети Хопфилда ограничена числом прототипов ($K \le 128$).
  \item Оценка на сбалансированном тесте завышает Precision относительно реального потока
        с долей fraud $<0.1\%$; для production необходима калибровка порога на несбалансированных данных.
\end{enumerate}

\textbf{Практическая значимость.}
Предложенный пайплайн применим как первый этап антифрод-фильтрации: MHSNet обучается на нормальных
операциях без полной разметки fraud, а гибрид с RF повышает итоговое качество при наличии
частично размеченных данных.

\section{Заключение}
\label{sec:conclusion}

В работе представлен воспроизводимый метод выявления мошеннических банковских транзакций
на основе гибридного пайплайна MHSNet (Kernel PCA $\rightarrow$ LOF $\rightarrow$ Hopfield)
и его комбинации с Random Forest. Эксперименты на реальном датасете CaixaBank подтверждают,
что MHSNet сопоставим с классическими unsupervised-методами, а гибридная модель достигает
$F_1 \approx 0.94$ и Recall@FPR$=5\% \approx 0.94$.

Направления дальнейших исследований: реализация SNN-ветки MHSNet, онлайн-обновление прототипов
при concept drift, интеграция графовых признаков и оценка в режиме, близком к real-time.

\section*{Доступность данных и кода}

Код, Jupyter notebook и инструкции воспроизведения доступны в репозитории проекта
\texttt{hopfield-network-for-fraud-nir}.
Датасет: \url{https://www.kaggle.com/datasets/computingvictor/transactions-fraud-datasets}.
Команда запуска основного эксперимента:
\texttt{python run.py --output outputs/caixabank}.

\section*{Благодарности}

Автор благодарит научного руководителя за поддержку при выполнении исследования.

\bibliographystyle{plainnat}
\bibliography{references}

\end{document}
